\documentclass[11pt,a4paper]{article}
\usepackage[utf8]{inputenc}
\usepackage[T1]{fontenc}
\usepackage[margin=2.5cm]{geometry}
\usepackage{amsmath}
\usepackage{enumitem}
\usepackage{hyperref}
\usepackage{parskip}

\hypersetup{colorlinks=true,linkcolor=blue,urlcolor=blue}

% allow long URLs to break at any character (no xurl available)
\makeatletter
\g@addto@macro\UrlBreaks{\do\a\do\b\do\c\do\d\do\e\do\f\do\g\do\h\do\i\do\j\do\k\do\l\do\m\do\n\do\o\do\p\do\q\do\r\do\s\do\t\do\u\do\v\do\w\do\x\do\y\do\z\do\A\do\B\do\C\do\D\do\E\do\F\do\G\do\H\do\I\do\J\do\K\do\L\do\M\do\N\do\O\do\P\do\Q\do\R\do\S\do\T\do\U\do\V\do\W\do\X\do\Y\do\Z}
\makeatother
\setlength{\emergencystretch}{2em}

\begin{document}

\begin{flushright}
Lucas Rover\\
Federal University of Technology --- Paraná (UTFPR)\\
Brazil\\

\bigskip
May, 2026
\end{flushright}

\bigskip

\noindent The Editors\\
\emph{Nature Communications}\\

\bigskip

Dear Dr.\ Bigorajski,

We submit our revised manuscript, ``\textbf{Same Prompt, Different Answer: Hidden Non-Determinism in LLM APIs Undermines Scientific Reproducibility}'', for further consideration at \emph{Nature Communications}. The revision incorporates 2{,}900 new controlled experiments across HumanEval (code generation), GSM8K (math reasoning), a cross-domain PubMed PM\textsubscript{2.5} probe, and a multi-turn extension to gpt-4o and DeepSeek; it reframes the unit of analysis from ``model'' to ``deployment stack'' as Reviewer 1 proposed; and it adds a per-field reproducibility analysis and an in-paper LLM-as-judge triangulation that together address every substantive concern raised in the first round. We thank you and the three reviewers for the constructive engagement and the clearly articulated path to a stronger manuscript.

This revision is substantial. We have addressed every point raised by Reviewers 1 and 3 (Reviewer 2 co-reviewed without independent comments) and incorporated the editor's specific request for experiments beyond summarisation. The eleven changes group into four substantive clusters:

\medskip
\noindent\textbf{A. Conceptual reframes and definitional clarifications.}
\begin{enumerate}[leftmargin=*]
  \item \textbf{Deployment stacks as the unit of analysis (R1.1).} Following Reviewer~1's central observation, we now treat the \emph{deployment stack} --- the tuple (model weights, provider, serving infrastructure, API layer) --- as the unit of analysis. This reframe is now explicit in the abstract, introduction, methods, tables, figure captions, and discussion. The Together AI quasi-isolation result is now positioned as the mechanistic core of the paper: it demonstrates that the stack, not the abstract model, is the carrier of variation.

  \item \textbf{Reframed quasi-isolation claim (R3.2).} The original strong claim that ``cloud deployment does not preclude reproducibility'' is reformulated as a defensible weaker claim, with explicit acknowledgement that the argument is partial evidence for closed-source proprietary stacks (which cannot be locally compared).

  \item \textbf{Cloud vs.\ production-serving infrastructure distinction (R1.4).} Methods now explicitly separates ``cloud deployment factors'' from ``production-serving infrastructure factors'', with concrete enumerations of each.
\end{enumerate}

\medskip
\noindent\textbf{B. Experimental expansion (2{,}900 new controlled experiments).}
\begin{enumerate}[leftmargin=*,resume]
  \item \textbf{Coding (HumanEval) and math reasoning (GSM8K) tasks (Editor; R1.8; R3.1).} Per the editor's specific request, we ran new controlled experiments on 30 HumanEval problems and 30 GSM8K problems, evaluated under our standard greedy-decoding protocol across all stacks. We report \emph{pass@1} for code and final-answer accuracy for math, alongside the same three-level reproducibility framework. The local--API gap persists in both new domains, with substantively analogous patterns to summarisation/extraction.

  \item \textbf{Multi-turn extension to gpt-4o and DeepSeek (R3.3).} We extended the three-turn refinement protocol to OpenAI gpt-4o and DeepSeek Chat, confirming that near-zero reproducibility under multi-turn regimes is universal across major cloud-served stacks rather than a Claude/Gemini-specific anomaly. RAG-extraction coverage remains as in the original analysis (Claude and Gemini); we acknowledge this scope explicitly in the Limitations.

  \item \textbf{Cross-domain PubMed PM\textsubscript{2.5}/respiratory-health probe (R3.4).} Ten additional abstracts demonstrate the reproducibility pattern persists outside the AI/ML domain (Claude Sonnet~4.5 EMR~=~0.010 [0.00, 0.03]).
\end{enumerate}

\medskip
\noindent\textbf{C. Validations, metrics, and concrete evidence.}
\begin{enumerate}[leftmargin=*,resume]
  \item \textbf{Field-level reproducibility analysis (R1.5).} We now compute every reproducibility metric per JSON output field. The result reveals an unexpected and stronger argument than we initially anticipated: BERTScore F1 is structurally saturated and unable to discriminate conclusion-relevant from metadata fields, while EMR exposes a paired Cohen's $d=+1.41$ gap between conclusion-relevant fields ($\overline{\text{EMR}}=0.455$) and metadata fields ($\overline{\text{EMR}}=0.684$). This resolves the apparent contradiction between ``textual not semantic'' and ``substantive content'': BERTScore alone is the wrong instrument; the three-level framework we propose is necessary.

  \item \textbf{Applied evidence-synthesis impact and two-judge LLM-as-judge triangulation (R1.9; R3.6).} We promote the PM\textsubscript{2.5}/respiratory-health analysis to a dedicated Results subsection (``Applied impact in evidence synthesis''). The substantive support is now \emph{self-contained within the present manuscript}: a two-judge blinded protocol using Claude Opus~4.7 (Anthropic) and gpt-4o (OpenAI) on a \emph{new} 30-case sample (distinct from the 23 study-level effects of the companion paper) against three pre-registered criteria. \textbf{Both judges classify the majority as truly contradictory} (Claude 22/30, Wilson 95\,\% CI 56--86\%; gpt-4o 27/30, 74--97\%; inter-judge Cohen's $\kappa$~=~0.29). The convergence of two independent provider families addresses the residual ``LLM-on-LLM'' concern more directly than either a single judge or a small single-modality human panel would. The companion paper (Rover \& de Souza Tadano, 2026, OSF) is cited only as larger-corpus context; the present manuscript's argument does not depend on it. A dedicated methodological extension applying multi-rater human validation against a domain-expert silver standard is planned as a follow-up paper.

  \item \textbf{Mechanism-by-stack mapping (R3.5).} A new Methods table (Table~3) maps each of the six known mechanisms of distributed-inference non-determinism to each deployment stack, with explicit labels (verifiable / likely / possible / mitigated) and an inferential caveat for closed-source stacks. A compact at-a-glance version of this mapping (Table~2) is included in the Results for in-text readability.

  \item \textbf{Concrete divergence examples (R1.12).} A new Box (Box~1, immediately after Fig.~5) shows side-by-side examples of cosmetic versus substantive divergences from the existing dataset.
\end{enumerate}

\medskip
\noindent\textbf{D. Editorial and reproducibility infrastructure.}
\begin{enumerate}[leftmargin=*,resume]
  \item \textbf{Submission checklists, ORCID, and data deposit.} We have completed the Code/Software submission checklist, the Machine Learning checklist, and the Reporting Summary update. All authors have linked ORCID identifiers on the Manuscript Tracking System. The full code and experimental records are deposited at \url{https://github.com/Roverlucas/genai-reproducibility-protocol} (MIT licence, release tag \texttt{v1.1}) and mirrored on Figshare with persistent DOI \href{https://doi.org/10.6084/m9.figshare.31653373}{10.6084/m9.figshare.31653373} (CC-BY 4.0). The GitHub repository holds the live source code; the Figshare deposit is the permanent, citable archive (code + all experimental records and source data) under that DOI, which becomes publicly accessible on acceptance. For peer-review access prior to public release, the private Figshare share URL is \href{https://figshare.com/s/3d17327cef1ae99ed37c}{https://figshare.com/s/3d17327cef1ae99ed37c}.
\end{enumerate}

The companion paper (Rover \& de Souza Tadano, 2026; OSF: \href{https://doi.org/10.17605/OSF.IO/VR934}{doi.org/10.17605/OSF.IO/VR934}) reports the full 500-abstract PM\textsubscript{2.5}/respiratory-health analysis. To avoid duplicate publication, the present manuscript reports only the aggregate propagation finding (23 effects appear/disappear depending on which run is used) as larger-corpus context, with full per-abstract data, silver-standard validation, and meta-analytic propagation residing in the companion manuscript, openly accessible at \href{https://doi.org/10.17605/OSF.IO/VR934}{doi.org/10.17605/OSF.IO/VR934}. The substantive evidence in the present manuscript is carried by the in-paper two-judge LLM-as-judge protocol on 30 new cases and the 10-abstract PubMed PM\textsubscript{2.5} probe; the companion paper is supporting context, not load-bearing.

The companion paper has not yet been submitted to any journal. Given the strong methodological complementarity between the two works---this one establishing the hidden non-determinism phenomenon and a provenance protocol that detects it; the companion paper quantifying its empirical impact on a 500-abstract evidence base in environmental health---\textbf{we would be very glad to submit the companion manuscript to \emph{Nature Communications} for parallel editorial consideration}, should you wish to evaluate the two as a coherent contribution.

Publishing the pair together at \emph{Nature Communications} would maximise the journal's impact in two concrete ways: \emph{(i)~the two works form a self-contained reproducibility story---methodological diagnosis (cause) and applied quantification (consequence)---and cross-citation between them within a single editorial venue amplifies both the methodological reach (this manuscript) and the policy-relevant impact (the companion paper) in ways neither paper achieves in isolation; (ii)~\emph{Nature Communications} has visibly led the conversation on reproducibility in computational science, and co-publishing the diagnosis with the applied-impact quantification would extend that leadership into LLM-based research, an area in which the field currently lacks an authoritative joint reference.} The companion paper has been openly preregistered at OSF (\href{https://doi.org/10.17605/OSF.IO/VR934}{doi.org/10.17605/OSF.IO/VR934}) for immediate editorial review, and we welcome the opportunity to position the pair within the same editorial venue.

A point-by-point response addressing every reviewer comment is provided as a separate document, together with the revised manuscript with all changes marked. Track changes were generated via \texttt{latexdiff} between the submitted version and the revised version. We hope the revisions meet the high standard you and the reviewers have set.

This manuscript has not been published elsewhere and is not under consideration by any other journal. All authors have approved the revision. We declare no competing interests.

\bigskip

\noindent Yours sincerely,

\bigskip\bigskip

\noindent \textbf{Lucas Rover}\\
on behalf of all authors\\
\href{mailto:lucasrover@alunos.utfpr.edu.br}{lucasrover@alunos.utfpr.edu.br}\\

\bigskip

\noindent \textit{On behalf of all authors:}\\
Lucas Rover (ORCID: \href{https://orcid.org/0000-0001-6641-9224}{0000-0001-6641-9224}),\\
Hugo Valadares Siqueira (ORCID: \href{https://orcid.org/0000-0002-1278-4602}{0000-0002-1278-4602}),\\
Eduardo Tadeu Bacalhau (ORCID: \href{https://orcid.org/0000-0002-3936-0375}{0000-0002-3936-0375}),\\
Anibal Tavares de Azevedo (ORCID: \href{https://orcid.org/0000-0003-1678-7795}{0000-0003-1678-7795}),\\
and Yara de Souza Tadano (ORCID: \href{https://orcid.org/0000-0002-3975-3419}{0000-0002-3975-3419})

\end{document}
